%%%%%%%%%%%%%%%%%%%%%%%%%%%%%%%%%%%%%%%%%%%%%%%%%%%%%%%%%%%%%%%%%
% GIÁO ÁN STEM: EPIDEMICSIM - MÔ PHỎNG LÂY LAN DỊCH BỆNH
% Môn: Toán 11 + Sinh học | Mô hình SIR và Hàm số Mũ
%%%%%%%%%%%%%%%%%%%%%%%%%%%%%%%%%%%%%%%%%%%%%%%%%%%%%%%%%%%%%%%%%
\documentclass[12pt,a4paper]{article}

% --- GÓI CƠ BẢN ---
\usepackage[utf8]{inputenc}
\usepackage[T1]{fontenc}
\usepackage[vietnamese]{babel}
\usepackage{graphicx}
\usepackage{amsmath,amssymb}
\usepackage{array}
\usepackage{booktabs}
\usepackage{multirow}
\usepackage{tabularx}
\usepackage{wrapfig}
\usepackage{ragged2e}
\usepackage{setspace}
\usepackage{tikz}
\usetikzlibrary{shapes,arrows,positioning,patterns,calc}
\usepackage{fontawesome5}
\usepackage{verbatim}
\usepackage{colortbl}

% --- FONT ---
\usepackage{newtxtext, newtxmath}
\usepackage[scaled=0.92]{helvet}
\usepackage{courier}
\usepackage{microtype}

% --- MÀU SẮC ---
\usepackage[svgnames]{xcolor}
\definecolor{primary}{RGB}{26,95,122}
\definecolor{secondary}{RGB}{21,152,149}
\definecolor{accent}{RGB}{255,107,53}
\definecolor{danger}{RGB}{194,55,40}
\definecolor{success}{RGB}{87,204,153}
\definecolor{warning}{RGB}{244,162,97}
\definecolor{TableHeader}{gray}{0.9}

\usepackage[left=2cm,right=2cm,top=2.5cm,bottom=2cm,headheight=15.6pt]{geometry}
\usepackage{fancyhdr}
\pagestyle{fancy}
\fancyhf{}
\renewcommand{\headrule}{\color{primary}\hrule width\textwidth height 1.5pt}
\fancyhead[L]{\color{primary}\footnotesize\sffamily\textbf{SỞ GD\&ĐT TP.HCM - HỘI THI STEM/STEAM CẤP THÀNH PHỐ}}
\fancyhead[R]{\color{primary}\footnotesize\sffamily\textbf{NĂM HỌC 2025-2026}}
\fancyfoot[C]{\sffamily\thepage}
\renewcommand{\footrule}{\color{primary}\hrule width\textwidth height 0.8pt}

% --- GÓI HỖ TRỢ ---
\usepackage{enumitem}
\usepackage{tcolorbox}
\tcbuselibrary{skins, breakable}
\usepackage[colorlinks=true,linkcolor=primary,urlcolor=accent,citecolor=success,
    pdftitle={Giao an STEM: EpidemicSim},pdfauthor={Nguyen Van Sang}]{hyperref}

% --- TIÊU ĐỀ ---
\usepackage{titlesec}
\titleformat{\subsection}{\normalfont\bfseries\color{accent}\sffamily}
{\llap{\color{accent}\rule{1.2em}{1.2em}\hspace{0.5em}}}{0em}{}
\titleformat{\subsubsection}{\normalfont\bfseries\itshape\color{secondary}\sffamily}
{\llap{\color{secondary}\rule{1em}{1em}\hspace{0.5em}}}{0em}{}

% --- MÔI TRƯỜNG ---
\def\phanI{TỔNG QUAN}
\def\phanII{CHUẨN BỊ CỦA GIÁO VIÊN}
\def\phanIII{KẾ HOẠCH DẠY HỌC ĐỀ XUẤT}
\def\phanIV{HƯỚNG DẪN HỌC SINH}
\def\phanV{CÁC PHỤ LỤC}

\newtcolorbox{competitionsection}[1]{
    enhanced, breakable, colback=white, colframe=primary, boxrule=0.8pt,
    fonttitle=\bfseries\sffamily\Huge, coltitle=white, colbacktitle=primary,
    attach boxed title to top left={xshift=1cm, yshift=-1mm-\tcboxedtitleheight/2},
    boxed title style={sharp corners, boxrule=0pt},
    title={PHẦN #1: \csname phan#1\endcsname},
    before upper={\addcontentsline{toc}{section}{Phần #1: \csname phan#1\endcsname}}
}

\newcolumntype{C}{>{\centering\arraybackslash}X}
\newcolumntype{L}{>{\raggedright\arraybackslash}X}
\newcommand{\header}[1]{\cellcolor{TableHeader}\sffamily\bfseries#1}

\newcommand{\teachingplan}[4]{
\noindent\begin{minipage}{0.12\textwidth}\centering\bfseries\sffamily\color{primary}#1\end{minipage}%
\begin{minipage}{0.23\textwidth}\sffamily\textbf{#2}\end{minipage}%
\begin{minipage}{0.35\textwidth}\textit{#3}\end{minipage}%
\begin{minipage}{0.25\textwidth}\small #4\end{minipage}
\vspace{1ex}\par\noindent\color{primary!30}\rule{\textwidth}{0.4pt}\vspace{1ex}}

\newtcolorbox{casestudy}{enhanced, breakable, colback=primary!5, colframe=primary!80, boxrule=1pt, arc=2mm,
    title=\sffamily\bfseries\color{accent}TÌNH HUỐNG THỰC TIỄN: ĐẠI DỊCH VÀ TOÁN HỌC, fonttitle=\bfseries\sffamily}
\newtcolorbox{keyformula}{enhanced, breakable, colback=accent!5, colframe=accent!80, boxrule=1pt, arc=2mm,
    title=\sffamily\bfseries\color{primary}KIẾN THỨC TRỌNG TÂM, fonttitle=\bfseries\sffamily}
\newtcolorbox{task}[1]{enhanced, breakable, colback=warning!10, colframe=warning!80, boxrule=1pt, arc=2mm,
    title=\sffamily\bfseries\color{secondary}NHIỆM VỤ: #1, fonttitle=\bfseries\sffamily}
\newtcolorbox{phieuhoc_tap_box}[1]{enhanced, breakable, colback=white, colframe=success, boxrule=1pt, arc=2mm,
    title=\sffamily\bfseries #1, fonttitle=\bfseries\sffamily, coltitle=white}

\setlist[itemize,1]{label=\textbullet, leftmargin=*, nosep}
\setlist[itemize,2]{label=--, leftmargin=*, nosep}
\setlist[description]{font=\bfseries\sffamily\color{primary}, nosep}

\begin{document}

% --- TRANG BÌA ---
\begin{titlepage}
\begin{tikzpicture}[remember picture, overlay]
\node[rectangle, fill=primary, anchor=north, minimum width=\paperwidth, minimum height=2cm] (header) at (current page.north){};
\node[anchor=south, yshift=2mm, text=white] at (header.south) {\Huge\sffamily\bfseries GIÁO ÁN DỰ THI STEM/STEAM CẤP THÀNH PHỐ};
\end{tikzpicture}

\vspace{2cm}
\begin{center}
{\LARGE\bfseries\color{primary}\sffamily SỞ GIÁO DỤC VÀ ĐÀO TẠO THÀNH PHỐ HỒ CHÍ MINH}\\
\vspace{0.5cm}
{\Large\bfseries\sffamily HỘI THI THIẾT KẾ CHỦ ĐỀ DẠY HỌC STEM/STEAM}\\
\vspace{0.3cm}
{\large\sffamily Năm học 2025 - 2026}

\vspace{2cm}

\begin{tcolorbox}[width=0.9\textwidth, colback=gray!5, colframe=primary, arc=2mm, boxrule=1pt]
    \centering
    {\LARGE\bfseries\sffamily CHỦ ĐỀ:} \\
    {\Huge\bfseries\color{danger}\sffamily EPIDEMICSIM}\\
    \vspace{0.25cm}
    {\Large\sffamily MÔ PHỎNG LÂY LAN DỊCH BỆNH}\\
    \vspace{0.25cm}
    {\large\sffamily Ứng dụng Mô hình SIR và Hàm số Mũ trong Dịch tễ học}
\end{tcolorbox}

\vspace{1.5cm}
{\Large\bfseries\sffamily MÔN: TOÁN + SINH HỌC - KHỐI: 11}

\vfill
\large
\begin{tabular}{rl}
\sffamily\bfseries Giáo viên: & Nguyễn Văn Sang\\
\sffamily\bfseries Đơn vị: & Trường THPT Nguyễn Hữu Cảnh\\
\sffamily\bfseries Số điện thoại: & 0389\,821\,115\\
\end{tabular}

\vspace{1cm}
{\large\sffamily Thành phố Hồ Chí Minh, tháng 11 năm 2025}
\end{center}
\end{titlepage}

\tableofcontents
\thispagestyle{empty}

%%%%% PHẦN I: TỔNG QUAN %%%%%
\begin{competitionsection}{I}
\vspace{0.5cm}
\subsection*{1.1. Thông tin cá nhân}
\begin{itemize}[leftmargin=*,noitemsep]
\item \textbf{Họ và tên:} Nguyễn Văn Sang
\item \textbf{Chuyên môn:} Toán học
\item \textbf{Thâm niên:} 15 năm giảng dạy THPT
\item \textbf{Đơn vị:} Trường THPT Nguyễn Hữu Cảnh, TP.HCM
\end{itemize}

\subsection*{1.2. Giới thiệu chủ đề}
\begin{itemize}[leftmargin=*,noitemsep]
\item \textbf{Tên chủ đề:} \textcolor{danger}{EpidemicSim: Mô phỏng lây lan dịch bệnh bằng Mô hình SIR và Hàm số Mũ}
\item \textbf{Bối cảnh thực tiễn:}
    \begin{itemize}
    \item Đại dịch COVID-19 cho thấy tầm quan trọng của việc hiểu cách dịch bệnh lây lan
    \item Các tin tức về ``làm phẳng đường cong'' và ``miễn dịch cộng đồng'' cần được giải thích bằng Toán học
    \item Học sinh cần thấy được ứng dụng thực tiễn của Hàm số Mũ trong y tế công cộng
    \end{itemize}
\item \textbf{Điểm mới và sáng tạo:}
    \begin{itemize}
    \item \textbf{Liên môn:} Kết hợp Toán học (Hàm số mũ, PT vi phân) + Sinh học (Dịch tễ học, Miễn dịch)
    \item \textbf{Mô phỏng động:} Công cụ web trực quan hóa quá trình lây lan theo thời gian thực
    \item \textbf{Ra quyết định:} Học sinh đóng vai chuyên gia y tế, đề xuất chính sách phòng dịch
    \item \textbf{Thời sự:} Liên hệ trực tiếp với kinh nghiệm COVID-19 của học sinh
    \end{itemize}
\end{itemize}

\subsection*{1.3. Căn cứ thiết kế}
\begin{minipage}{0.48\textwidth}
\textbf{Căn cứ pháp lý:}
\begin{itemize}[leftmargin=*,noitemsep]
\item Chương trình GDPT 2018 môn Toán lớp 11 (Hàm số mũ và logarit)
\item Công văn 5555/BGDĐT-GDTrH về dạy học STEM
\item Kế hoạch số 4500/KH-SGDĐT của Sở GD\&ĐT TP.HCM
\end{itemize}
\end{minipage}
\hfill
\begin{minipage}{0.48\textwidth}
\textbf{Cơ sở lý luận:}
\begin{itemize}[leftmargin=*,noitemsep]
\item Mô hình SIR (Kermack-McKendrick, 1927) trong Dịch tễ học
\item Mô hình 5E (Engage, Explore, Explain, Elaborate, Evaluate)
\item Học tập dựa trên mô phỏng (Simulation-Based Learning)
\end{itemize}
\end{minipage}

\subsection*{1.4. Mục tiêu giáo dục}
\begin{center}
\begin{tabularx}{\linewidth}{@{} L L L @{}}
\toprule
\header{Kiến thức} & \header{Kỹ năng} & \header{Thái độ} \\
\midrule
\parbox[t]{\dimexpr1.\linewidth-2\tabcolsep}{
    \begin{itemize}[leftmargin=*,noitemsep,topsep=1ex]
        \item Hiểu mô hình SIR và các tham số dịch tễ
        \item Áp dụng hàm số mũ để mô tả tăng trưởng dịch
        \item Tính hệ số lây nhiễm $R_0$ và ngưỡng miễn dịch cộng đồng
    \end{itemize}
} &
\parbox[t]{\dimexpr1.\linewidth-2\tabcolsep}{
    \begin{itemize}[leftmargin=*,noitemsep,topsep=1ex]
        \item Mô hình hóa bài toán thực tế
        \item Sử dụng công nghệ mô phỏng
        \item Đọc hiểu đồ thị, phân tích kịch bản
        \item Ra quyết định dựa trên dữ liệu
    \end{itemize}
} &
\parbox[t]{\dimexpr1.\linewidth-2\tabcolsep}{
    \begin{itemize}[leftmargin=*,noitemsep,topsep=1ex]
        \item Ý thức phòng chống dịch bệnh
        \item Tư duy phản biện với thông tin y tế
        \item Trách nhiệm với sức khỏe cộng đồng
    \end{itemize}
} \\
\bottomrule
\end{tabularx}
\end{center}

\subsection*{1.5. Cấu trúc chương trình}
\begin{center}
\begin{tikzpicture}[node distance=1.5cm, every node/.style={font=\sffamily}]
\node (start) [rectangle, rounded corners, fill=primary!20, text width=4.5cm, align=center, minimum height=1.2cm] {\textbf{Khám phá thực tiễn}\\ Đại dịch COVID-19 đã lây lan như thế nào?};
\node (mid1) [rectangle, rounded corners, fill=accent!20, text width=4.5cm, align=center, below of=start, yshift=-1.2cm, minimum height=1.2cm] {\textbf{Xây dựng mô hình Toán}\\ Hàm số mũ, Mô hình SIR};
\node (mid2) [rectangle, rounded corners, fill=primary!20, text width=4.5cm, align=center, below of=mid1, yshift=-1.2cm, minimum height=1.2cm] {\textbf{Mô phỏng và phân tích}\\ Ứng dụng EpidemicSim.html};
\node (end) [rectangle, rounded corners, fill=accent!20, text width=4.5cm, align=center, below of=mid2, yshift=-1.2cm, minimum height=1.2cm] {\textbf{Ra quyết định}\\ Đề xuất chính sách phòng dịch};
\draw[->, thick, color=primary] (start) -- (mid1) node[midway, right, text width=3cm, align=center] {\footnotesize \textbf{S} (Sinh học)};
\draw[->, thick, color=primary] (mid1) -- (mid2) node[midway, right, text width=3cm, align=center] {\footnotesize \textbf{T+M} (Công nghệ + Toán)};
\draw[->, thick, color=primary] (mid2) -- (end) node[midway, right, text width=3cm, align=center] {\footnotesize \textbf{E} (Kỹ thuật/Chính sách)};
\end{tikzpicture}
\end{center}
\end{competitionsection}
\newpage

%%%%% PHẦN II: CHUẨN BỊ %%%%%
\begin{competitionsection}{II}
\vspace{0.5cm}
\subsection*{2.1. Thiết bị và công nghệ}
\begin{itemize}[leftmargin=*,noitemsep]
\item \textbf{Phần cứng:}
    \begin{itemize}
    \item Máy tính hoặc điện thoại thông minh có trình duyệt web
    \item Máy chiếu tương tác (nếu có)
    \end{itemize}
\item \textbf{Phần mềm và nền tảng:}
    \begin{itemize}
    \item \textbf{Học liệu số chính:} Ứng dụng web \texttt{EpidemicSim.html} hoặc \href{https://stem-1h8.pages.dev/EpidemicSim.html}{stem-1h8.pages.dev/EpidemicSim.html}
    \item \textit{Yêu cầu:} Trình duyệt web (Chrome, Firefox, Edge...)
    \item Padlet/Jamboard để chia sẻ kết quả nhóm
    \end{itemize}
\item \textbf{Tài liệu:}
    \begin{itemize}
    \item Phiếu học tập số 1: Khám phá Mô hình SIR
    \item Phiếu học tập số 2: Hướng dẫn sử dụng EpidemicSim
    \item Phiếu khảo sát trước/sau học
    \item Rubric đánh giá sản phẩm
    \end{itemize}
\end{itemize}

\subsection*{2.2. Kịch bản giảng dạy (2 tiết = 90 phút)}
\begin{center}
\begin{tabularx}{\linewidth}{@{} L L L L L @{}}
\toprule
\header{Thời lượng} & \header{Hoạt động} & \header{Phương pháp} & \header{Công cụ} & \header{Đánh giá} \\
\midrule
10' & Khởi động: Dịch COVID-19 & Đàm thoại & Clip + Số liệu thực & Quan sát hứng thú \\
15' & Khám phá Mô hình SIR & Khám phá có hướng dẫn & PHT số 1 & Phiếu trả lời \\
\rowcolor{primary!10}
10' & Giới thiệu EpidemicSim & Demo ứng dụng & Web App (3 tabs) & Câu hỏi phản hồi \\
\rowcolor{primary!10}
25' & Mô phỏng và phân tích & Học qua dự án & EpidemicSim & Nhật ký học tập \\
\rowcolor{primary!10}
15' & So sánh kịch bản & Hợp tác nhóm & Tab 2: So sánh & Peer assessment \\
10' & Trình bày sản phẩm & Thuyết trình & Google Slides & Rubric đánh giá \\
5' & Tổng kết & Suy ngẫm & Tab 3: Học tập & Khảo sát sau học \\
\bottomrule
\end{tabularx}
\end{center}

\subsection*{2.3. Đánh giá đa chiều}
\begin{minipage}[t]{0.48\textwidth}
\textbf{Đánh giá quá trình:}
\begin{itemize}[leftmargin=*,noitemsep]
\item \textbf{Phiếu quan sát nhóm:} Mức độ hợp tác, tư duy phản biện
\item \textbf{Nhật ký học tập:} Suy nghĩ cá nhân, câu hỏi đặt ra
\item \textbf{Đánh giá đồng cấp:} Học sinh đánh giá lẫn nhau
\end{itemize}
\end{minipage}
\hfill
\begin{minipage}[t]{0.48\textwidth}
\textbf{Đánh giá sản phẩm:}
\begin{itemize}[leftmargin=*,noitemsep]
\item \textbf{Rubric 4 tiêu chí:}
    \begin{itemize}
    \item Hiểu mô hình SIR (4 điểm)
    \item Sử dụng công nghệ mô phỏng (4 điểm)
    \item Phân tích và ra quyết định (4 điểm)
    \item Trình bày và bảo vệ (4 điểm)
    \end{itemize}
\end{itemize}
\end{minipage}

\subsection*{2.4. Kế hoạch dự phòng}
\begin{itemize}[leftmargin=*,noitemsep]
\item \textbf{Tình huống 1:} Mất điện hoặc mất kết nối Internet
    \begin{itemize}
    \item Ứng dụng \texttt{EpidemicSim.html} \textbf{không cần Internet}, chỉ cần mở bằng trình duyệt.
    \item Dự phòng: Sử dụng bảng tính Excel với công thức SIR đã chuẩn bị sẵn.
    \end{itemize}
\item \textbf{Tình huống 2:} Học sinh không có thiết bị cá nhân
    \begin{itemize}
    \item Sắp xếp nhóm 3-4 học sinh/1 máy tính
    \item Luân phiên sử dụng thiết bị
    \end{itemize}
\item \textbf{Tình huống 3:} Học sinh tiếp thu nhanh
    \begin{itemize}
    \item Giao nhiệm vụ mở rộng: Nghiên cứu mô hình SEIR (thêm nhóm Exposed)
    \item Tìm hiểu số liệu thực từ WHO
    \end{itemize}
\end{itemize}
\end{competitionsection}

%%%%% PHẦN III: KẾ HOẠCH DẠY HỌC %%%%%
\begin{competitionsection}{III}
\vspace{0.5cm}
\subsection*{3.1. Tiến trình hai tiết học (Tóm tắt)}
\teachingplan{10'}{Khởi động}{Video/Số liệu COVID-19: ``Tại sao số ca tăng theo cấp số nhân?''}{Công nghệ: Clip + Infographic}
\teachingplan{15'}{Khám phá mô hình SIR}{Phiếu học tập: S-I-R là gì? Ý nghĩa các tham số}{PP: Khám phá có hướng dẫn}
\teachingplan{10'}{Demo EpidemicSim}{Giới thiệu 3 tabs: Mô phỏng, So sánh, Học tập}{Công nghệ: Web App}
\teachingplan{25'}{Mô phỏng và phân tích}{Thực hành mô phỏng, thay đổi tham số, quan sát kết quả}{PP: Học qua dự án}
\teachingplan{15'}{So sánh kịch bản}{So sánh: Không can thiệp vs Giãn cách vs Tiêm chủng}{Công cụ: Biểu đồ so sánh}
\teachingplan{10'}{Trình bày}{Thuyết trình: Đề xuất chính sách phòng dịch cho địa phương}{PP: Thuyết trình nhóm}
\teachingplan{5'}{Tổng kết}{Kiến thức SIR, Hàm mũ, Ý nghĩa thực tiễn}{Tab 3: Học tập}

\subsection*{3.2. Kịch bản chi tiết}

\begin{casestudy}
Năm 2020, đại dịch COVID-19 bùng phát và lan rộng toàn cầu. Tại Việt Nam, từ 0 ca ban đầu, số ca nhiễm tăng theo cấp số nhân, đặc biệt trong làn sóng Delta.

\textbf{Thách thức:} Là một chuyên gia dịch tễ học, bạn cần:
\begin{itemize}[leftmargin=*,noitemsep]
\item Hiểu cách dịch bệnh lây lan theo mô hình toán học
\item Dự đoán ``đỉnh dịch'' sẽ xảy ra khi nào
\item Đề xuất biện pháp can thiệp (giãn cách, tiêm chủng) để ``làm phẳng đường cong''
\end{itemize}
\end{casestudy}

\subsubsection*{Hoạt động 1: Khởi động - Kết nối thực tiễn (10 phút)}
\begin{itemize}[leftmargin=*,noitemsep]
\item \textbf{Bước 1:} Giáo viên trình chiếu biểu đồ số ca COVID-19 tại Việt Nam (2020-2022)
\item \textbf{Bước 2:} Đặt câu hỏi kích thích tư duy:
    \begin{itemize}
    \item \textit{``Tại sao số ca nhiễm tăng rất nhanh ở giai đoạn đầu?''}
    \item \textit{``Miễn dịch cộng đồng là gì? Cần bao nhiêu \% dân số có miễn dịch?''}
    \end{itemize}
\item \textbf{Bước 3:} Học sinh thảo luận cặp đôi và chia sẻ ý kiến
\item \textbf{Chuyển tiếp:} \textit{``Toán học có thể giúp chúng ta dự đoán và kiểm soát dịch bệnh!''}
\end{itemize}

\subsubsection*{Hoạt động 2: Khám phá Mô hình SIR (15 phút)}
\begin{task}{Nhóm}
\textbf{Nhiệm vụ:} Tìm hiểu mô hình SIR qua Phiếu học tập số 1

\textbf{Nội dung:}
\begin{itemize}[leftmargin=*,noitemsep]
\item \textbf{S} (Susceptible): Người có thể bị lây nhiễm
\item \textbf{I} (Infected): Người đang nhiễm bệnh
\item \textbf{R} (Recovered): Người đã khỏi bệnh và có miễn dịch
\end{itemize}
\end{task}

\begin{keyformula}
\textbf{Hệ phương trình vi phân SIR:}
$$\frac{dS}{dt} = -\beta \cdot \frac{S \cdot I}{N}, \quad \frac{dI}{dt} = \beta \cdot \frac{S \cdot I}{N} - \gamma \cdot I, \quad \frac{dR}{dt} = \gamma \cdot I$$

\textbf{Các tham số:}
\begin{itemize}[leftmargin=*,noitemsep]
\item $\beta$ = Tỷ lệ lây nhiễm (số người bị lây/người nhiễm/ngày)
\item $\gamma$ = Tỷ lệ hồi phục (1/thời gian nhiễm trung bình)
\item $R_0 = \beta/\gamma$ = Hệ số lây nhiễm cơ bản
\item Nếu $R_0 > 1$: Dịch bùng phát | Nếu $R_0 < 1$: Dịch tự tắt
\end{itemize}

\textbf{Giai đoạn đầu (Hàm số mũ):}
$$I(t) = I_0 \cdot e^{(\beta - \gamma) \cdot t} = I_0 \cdot e^{r \cdot t}$$
\end{keyformula}

\subsubsection*{Hoạt động 3: Demo ứng dụng EpidemicSim (10 phút)}
\begin{itemize}[leftmargin=*,noitemsep]
\item \textbf{Bước 1:} Giáo viên demo ứng dụng web EpidemicSim.html
\item \textbf{Bước 2:} Giới thiệu 3 tab:
    \begin{itemize}
    \item \textbf{Tab 1 - Mô phỏng:} Nhập tham số, xem biểu đồ SIR động
    \item \textbf{Tab 2 - So sánh kịch bản:} Thêm kịch bản, so sánh đỉnh dịch
    \item \textbf{Tab 3 - Học tập:} Công thức, kiến thức dịch tễ học
    \end{itemize}
\item \textbf{Bước 3:} Thay đổi tham số $\beta$ và quan sát sự thay đổi của đồ thị
\end{itemize}

\subsubsection*{Hoạt động 4: Mô phỏng và phân tích (25 phút)}
\begin{task}{Nhóm}
\textbf{Nhiệm vụ:} Sử dụng EpidemicSim để mô phỏng 3 kịch bản

\textbf{Kịch bản 1 - Không can thiệp:} $\beta = 0.4$, $\gamma = 0.1$, Tiêm chủng = 0\%

\textbf{Kịch bản 2 - Giãn cách xã hội:} $\beta = 0.2$, $\gamma = 0.1$, Tiêm chủng = 0\%

\textbf{Kịch bản 3 - Tiêm chủng 50\%:} $\beta = 0.4$, $\gamma = 0.1$, Tiêm chủng = 50\%

\textbf{Yêu cầu:} Ghi lại đỉnh dịch (số ca max) và ngày đỉnh dịch của mỗi kịch bản.
\end{task}

\subsubsection*{Hoạt động 5: So sánh và trình bày (15 phút)}
\begin{itemize}[leftmargin=*,noitemsep]
\item \textbf{Bước 1:} 2 nhóm trình bày kết quả so sánh 3 kịch bản
\item \textbf{Bước 2:} Các nhóm khác đặt câu hỏi phản biện
\item \textbf{Bước 3:} Giáo viên đặt câu hỏi mở rộng:
    \begin{itemize}
    \item \textit{``Tại sao giãn cách xã hội làm giảm $\beta$?''}
    \item \textit{``Cần tiêm chủng bao nhiêu \% để đạt miễn dịch cộng đồng?''}
    \end{itemize}
\end{itemize}

\subsubsection*{Hoạt động 6: Tổng kết (5 phút)}
\begin{itemize}[leftmargin=*,noitemsep]
\item \textbf{Kiến thức trọng tâm:}
    \begin{itemize}
    \item Mô hình SIR giúp dự đoán diễn biến dịch bệnh
    \item Hàm số mũ mô tả giai đoạn đầu của dịch
    \item $R_0 > 1$ là điều kiện để dịch bùng phát
    \item Giãn cách và tiêm chủng là hai biện pháp ``làm phẳng đường cong''
    \end{itemize}
\item \textbf{Bài tập về nhà:}
    \begin{itemize}
    \item Tìm hiểu $R_0$ của các bệnh khác (Sởi, Cúm, Ebola)
    \item Tính ngưỡng miễn dịch cộng đồng $H = 1 - 1/R_0$ cho mỗi bệnh
    \end{itemize}
\end{itemize}
\end{competitionsection}
\newpage

%%%%% PHẦN IV: HƯỚNG DẪN HỌC SINH %%%%%
\begin{competitionsection}{IV}
\vspace{0.5cm}
\subsection*{4.1. Hướng dẫn chi tiết cho học sinh}
\subsubsection*{Giai đoạn 1: Tìm hiểu Mô hình SIR}
\begin{itemize}[leftmargin=*,noitemsep]
\item \textbf{Bước 1:} Đọc kỹ tình huống đại dịch COVID-19
\item \textbf{Bước 2:} Xác định 3 nhóm dân số:
    \begin{itemize}
    \item \textbf{S (Susceptible):} Người có thể bị lây (chưa có miễn dịch)
    \item \textbf{I (Infected):} Người đang nhiễm bệnh
    \item \textbf{R (Recovered):} Người đã khỏi bệnh và có miễn dịch
    \end{itemize}
\item \textbf{Bước 3:} Viết công thức $R_0 = \beta / \gamma$ và giải thích ý nghĩa
\item \textbf{Bước 4:} Tính ngưỡng miễn dịch cộng đồng $H = 1 - 1/R_0$
\end{itemize}

\subsubsection*{Giai đoạn 2: Sử dụng EpidemicSim}
\begin{itemize}[leftmargin=*,noitemsep]
\item \textbf{Bước 1: Mở ứng dụng}
    \begin{itemize}
    \item Mở file \texttt{EpidemicSim.html} bằng trình duyệt hoặc truy cập website
    \end{itemize}
\item \textbf{Bước 2: Tab 1 - Mô phỏng}
    \begin{itemize}
    \item Nhập dân số (N = 10,000)
    \item Nhập số ca ban đầu ($I_0$ = 10)
    \item Điều chỉnh $\beta$ và $\gamma$ bằng thanh trượt
    \item Nhấn ``Chạy mô phỏng'' và quan sát biểu đồ
    \end{itemize}
\item \textbf{Bước 3: Tab 2 - So sánh kịch bản}
    \begin{itemize}
    \item Thêm các kịch bản khác nhau (thay đổi $\beta$, tỷ lệ tiêm chủng)
    \item Nhấn ``So sánh'' để xem biểu đồ cột so sánh đỉnh dịch
    \end{itemize}
\item \textbf{Bước 4: Tab 3 - Học tập}
    \begin{itemize}
    \item Đọc công thức SIR và Hàm số mũ
    \item Tìm hiểu ngưỡng miễn dịch cộng đồng của các bệnh khác nhau
    \end{itemize}
\end{itemize}

\subsubsection*{Giai đoạn 3: Phân tích và trình bày}
\begin{itemize}[leftmargin=*,noitemsep]
\item \textbf{Bước 1:} So sánh các kịch bản dựa trên:
    \begin{itemize}
    \item Đỉnh dịch (số ca nhiễm cao nhất)
    \item Ngày đỉnh dịch
    \item Tổng số ca nhiễm
    \end{itemize}
\item \textbf{Bước 2:} Đề xuất chính sách phòng dịch cho địa phương
\item \textbf{Bước 3:} Chuẩn bị trình bày 3-5 phút
\end{itemize}

\subsection*{4.2. Gợi ý mở rộng}
\begin{itemize}[leftmargin=*,noitemsep]
\item \textbf{Tìm hiểu sâu về Dịch tễ học:}
    \begin{itemize}
    \item Mô hình SEIR (thêm nhóm Exposed - đã nhiễm nhưng chưa lây)
    \item Số liệu thực từ WHO và Bộ Y tế Việt Nam
    \end{itemize}
\item \textbf{Liên hệ môn Sinh học:}
    \begin{itemize}
    \item Cơ chế lây lan virus qua giọt bắn, tiếp xúc
    \item Nguyên lý hoạt động của vaccine
    \item Miễn dịch tự nhiên vs miễn dịch nhân tạo
    \end{itemize}
\item \textbf{Ứng dụng thực tế:}
    \begin{itemize}
    \item Theo dõi tin tức dịch bệnh và phân tích bằng kiến thức đã học
    \item Tham gia các hoạt động tuyên truyền phòng dịch tại trường
    \end{itemize}
\end{itemize}
\end{competitionsection}
\newpage

%%%%% PHẦN V: PHỤ LỤC %%%%%
\begin{competitionsection}{V}
\vspace{0.5cm}
\subsection*{Phụ lục 1: Phiếu học tập số 1 - Khám phá Mô hình SIR}
\begin{phieuhoc_tap_box}{Phụ lục 1: Phiếu học tập số 1 - Mô hình SIR}
\textbf{Tình huống:} Một thành phố có 10,000 dân. Có 10 ca nhiễm bệnh X ban đầu.

\vspace{0.5cm}
\textbf{Bài 1: Xác định 3 nhóm dân số}
\begin{itemize}[leftmargin=*,noitemsep]
\item \textbf{S} (Susceptible) = \underline{\hspace{6cm}} người
\item \textbf{I} (Infected) = \underline{\hspace{6cm}} người
\item \textbf{R} (Recovered) = \underline{\hspace{6cm}} người
\item Kiểm tra: S + I + R = \underline{\hspace{2cm}} = N (dân số)
\end{itemize}

\textbf{Bài 2: Tính $R_0$ (điền vào chỗ trống)}
\begin{itemize}[leftmargin=*,noitemsep]
\item Biết $\beta = 0.3$ (tỷ lệ lây nhiễm/ngày), $\gamma = 0.1$ (tỷ lệ hồi phục/ngày)
\item $R_0 = \beta / \gamma = \underline{\hspace{2cm}} / \underline{\hspace{2cm}} = \underline{\hspace{2cm}}$
\item Ý nghĩa: Mỗi người nhiễm trung bình lây cho \underline{\hspace{2cm}} người khác
\end{itemize}

\textbf{Bài 3: Ngưỡng miễn dịch cộng đồng}
\begin{itemize}[leftmargin=*,noitemsep]
\item Công thức: $H = 1 - 1/R_0 = 1 - 1/\underline{\hspace{1cm}} = \underline{\hspace{2cm}}$
\item Ý nghĩa: Cần \underline{\hspace{2cm}}\% dân số có miễn dịch để dịch tự tắt
\end{itemize}

\textbf{Bài 4: Dự đoán}
\begin{itemize}[leftmargin=*,noitemsep]
\item Dịch có bùng phát không? ($R_0 > 1$ hay $R_0 < 1$): \underline{\hspace{4cm}}
\item Nếu không can thiệp, dự đoán số ca nhiễm sau 30 ngày: \underline{\hspace{4cm}}
\end{itemize}
\end{phieuhoc_tap_box}

\newpage
\subsection*{Phụ lục 2: Phiếu học tập số 2 - Hướng dẫn sử dụng EpidemicSim}
\begin{phieuhoc_tap_box}{Phụ lục 2: Nhiệm vụ với Web App EpidemicSim.html}
\textbf{I. Giới thiệu ứng dụng}
\begin{itemize}[leftmargin=*,noitemsep]
\item \textbf{Địa chỉ:} \href{https://stem-1h8.pages.dev/EpidemicSim.html}{stem-1h8.pages.dev/EpidemicSim.html} hoặc file \texttt{EpidemicSim.html}
\item \textbf{Cấu trúc 3 Tab:} (1) Mô phỏng, (2) So sánh kịch bản, (3) Học tập
\end{itemize}

\textbf{II. Nhiệm vụ thực hành}
\begin{itemize}[leftmargin=*,noitemsep]
\item \textbf{Bước 1: Tab ``Mô phỏng''}
    \begin{itemize}
    \item Nhập: N = 10,000 | $I_0$ = 10 | $\beta$ = 0.3 | $\gamma$ = 0.1 | Tiêm chủng = 0\%
    \item Nhấn ``Chạy mô phỏng''
    \item Ghi lại: Đỉnh dịch = \underline{\hspace{2cm}} ca | Ngày đỉnh = \underline{\hspace{2cm}}
    \end{itemize}
\item \textbf{Bước 2: Tab ``So sánh kịch bản''}
    \begin{itemize}
    \item Thêm Kịch bản 2: $\beta$ = 0.2 (giãn cách xã hội)
    \item Thêm Kịch bản 3: Tiêm chủng = 50\%
    \item Nhấn ``So sánh''
    \item Kịch bản nào tốt nhất? \underline{\hspace{6cm}}
    \end{itemize}
\item \textbf{Bước 3: Tab ``Học tập''}
    \begin{itemize}
    \item Đọc công thức $H = 1 - 1/R_0$
    \item Với COVID-19 ($R_0 \approx 2.5$), ngưỡng miễn dịch = \underline{\hspace{2cm}}\%
    \end{itemize}
\end{itemize}

\textbf{III. Kết luận}

Chính sách phòng dịch nhóm em đề xuất: \underline{\hspace{0.8\textwidth}}

Lý do: \underline{\hspace{0.9\textwidth}}
\end{phieuhoc_tap_box}

\newpage
\subsection*{Phụ lục 3: Rubric đánh giá sản phẩm}
\begin{tabularx}{\linewidth}{@{} >{\raggedright\arraybackslash}p{3cm} >{\raggedright\arraybackslash}X >{\raggedright\arraybackslash}X >{\raggedright\arraybackslash}X @{}}
    \toprule
    \rowcolor{TableHeader}
    \sffamily\bfseries Tiêu chí & \sffamily\bfseries Xuất sắc (4đ) & \sffamily\bfseries Khá (3đ) & \sffamily\bfseries Cần cải thiện (1-2đ) \\
    \midrule
    
    \textbf{Hiểu mô hình SIR} 
    & 
    \begin{itemize}[leftmargin=*,noitemsep,topsep=1ex]
        \item Giải thích đúng S, I, R
        \item Tính đúng $R_0$ và $H$
        \item Hiểu ý nghĩa các tham số
    \end{itemize}
    & 
    \begin{itemize}[leftmargin=*,noitemsep,topsep=1ex]
        \item Hiểu cơ bản S, I, R
        \item Tính đúng công thức
    \end{itemize}
    & 
    \begin{itemize}[leftmargin=*,noitemsep,topsep=1ex]
        \item Nhầm lẫn các khái niệm
        \item Tính sai công thức
    \end{itemize} \\
    \addlinespace
    
    \textbf{Sử dụng EpidemicSim} 
    & 
    \begin{itemize}[leftmargin=*,noitemsep,topsep=1ex]
        \item Thành thạo cả 3 tabs
        \item So sánh 3+ kịch bản
    \end{itemize}
    & 
    \begin{itemize}[leftmargin=*,noitemsep,topsep=1ex]
        \item Sử dụng tốt Tab 1, 2
        \item So sánh 2 kịch bản
    \end{itemize}
    & 
    \begin{itemize}[leftmargin=*,noitemsep,topsep=1ex]
        \item Chỉ dùng được Tab 1
        \item Không so sánh được
    \end{itemize} \\
    \addlinespace
    
    \textbf{Phân tích \& Ra quyết định} 
    & 
    \begin{itemize}[leftmargin=*,noitemsep,topsep=1ex]
        \item Đề xuất chính sách hợp lý
        \item Lý giải dựa trên dữ liệu
    \end{itemize}
    & 
    \begin{itemize}[leftmargin=*,noitemsep,topsep=1ex]
        \item Đề xuất cơ bản
        \item Có lý giải
    \end{itemize}
    & 
    \begin{itemize}[leftmargin=*,noitemsep,topsep=1ex]
        \item Không đề xuất được
        \item Thiếu lý giải
    \end{itemize} \\
    \addlinespace
    
    \textbf{Trình bày} 
    & 
    \begin{itemize}[leftmargin=*,noitemsep,topsep=1ex]
        \item Mạch lạc, logic
        \item Có biểu đồ minh họa
    \end{itemize}
    & 
    \begin{itemize}[leftmargin=*,noitemsep,topsep=1ex]
        \item Trình bày rõ ràng
        \item Có hỗ trợ trực quan
    \end{itemize}
    & 
    \begin{itemize}[leftmargin=*,noitemsep,topsep=1ex]
        \item Trình bày rời rạc
        \item Không có minh họa
    \end{itemize} \\
    \bottomrule
\end{tabularx}
\vspace{1em}
\noindent\textbf{Tổng điểm:} \underline{\hspace{2cm}}/16 điểm

\newpage
\subsection*{Phụ lục 4: Bảng tham khảo $R_0$ các bệnh thường gặp}
\begin{center}
\begin{tabularx}{0.95\linewidth}{@{} l C C C @{}}
\toprule
\header{Bệnh} & \header{$R_0$ (ước tính)} & \header{Ngưỡng miễn dịch $H$} & \header{Đặc điểm} \\
\midrule
COVID-19 (gốc) & 2.5 - 3.0 & 60-67\% & Lây qua giọt bắn, tiếp xúc \\
COVID-19 (Delta) & 5 - 8 & 80-88\% & Biến thể lây lan nhanh hơn \\
Cúm mùa & 1.2 - 1.5 & 17-33\% & Lây theo mùa \\
Sởi & 12 - 18 & 92-94\% & Rất dễ lây, cần tiêm chủng cao \\
Ebola & 1.5 - 2.5 & 33-60\% & Lây qua dịch cơ thể \\
HIV/AIDS & 2 - 5 & 50-80\% & Lây qua máu, tình dục \\
\bottomrule
\end{tabularx}
\end{center}

\subsection*{Phụ lục 5: Mã nguồn thuật toán SIR (JavaScript)}
\begin{verbatim}
function simulateSIR(N, I0, beta, gamma, days) {
    let S = N - I0;
    let I = I0;
    let R = 0;
    
    const data = { days: [0], S: [S], I: [I], R: [R] };
    
    for (let t = 1; t <= days; t++) {
        const dS = -beta * S * I / N;
        const dI = beta * S * I / N - gamma * I;
        const dR = gamma * I;
        
        S = Math.max(0, S + dS);
        I = Math.max(0, I + dI);
        R = Math.min(N, R + dR);
        
        data.days.push(t);
        data.S.push(Math.round(S));
        data.I.push(Math.round(I));
        data.R.push(Math.round(R));
    }
    return data;
}
\end{verbatim}

\textbf{Hướng dẫn:} Học sinh có thể xem mã nguồn đầy đủ bằng cách nhấn F12 (DevTools) hoặc View Source.
\end{competitionsection}

\end{document}
